% =============================================================================
% Supporting Information
% Coordinated Vibronic Light-Harvesting & Polaritonic Interface
% for Symbiotic Quantum Agrivoltaics
% Target: Nature Energy
% Date:    June 18, 2026 (First Draft)
% Authors: Teguia Kouam S.C., Goumai Vedekoi T., Tchapet Njafa J.-P.,
%          Nguenang J.-P., Nana Engo S.G.
% =============================================================================

\documentclass[11pt,a4paper]{article}

% ---- Required packages ----
\usepackage[utf8]{inputenc}
\usepackage[T1]{fontenc}
\usepackage{amsmath,amssymb,amsfonts,graphicx,booktabs,physics,bm,multirow}
\usepackage{siunitx}
\DeclareSIUnit\year{yr}
\usepackage{xcolor}
\usepackage[margin=2.5cm]{geometry}
\usepackage[colorlinks=true,linkcolor=blue,citecolor=blue,urlcolor=blue]{hyperref}
\usepackage[capitalise,nameinlink]{cleveref}
\usepackage{authblk}
\usepackage{xr}

% ---- Cross-reference to main manuscript ----
\externaldocument{Manuscript}

% ---- siunitx v3/v4 configuration ----
\sisetup{
    locale = US,
    detect-all = true,
    group-minimum-digits = 4,
    separate-uncertainty = true,
    multi-part-units = single,
    range-phrase = {--},
    range-units = single,
    quantity-product = \,,
}

% ---- SI numbering conventions ----
\renewcommand{\theequation}{S\arabic{equation}}
\renewcommand{\thefigure}{S\arabic{figure}}
\renewcommand{\thetable}{S\arabic{table}}
\renewcommand{\thesection}{S\arabic{section}}

% =============================================================================
% DOCUMENT
% =============================================================================

\begin{document}

\title{\textbf{Supporting Information}\\[0.5em]
\large Coordinated Vibronic Light-Harvesting and Polaritonic Interface for Symbiotic Quantum Agrivoltaics}

\author{Steve Cabrel Teguia Kouam}
\affil{Department of Physics, Faculty of Science, University of Douala, Cameroon}
\author{Theodore Goumai Vedekoi}
\affil{Department of Physics, Faculty of Science, University of Yaound\'{e}~I, Cameroon}
\author{Jean-Pierre Tchapet Njafa}
\affil{Department of Physics, Faculty of Science, University of Yaound\'{e}~I, Cameroon}
\author{Jean-Pierre Nguenang}
\affil{Department of Physics, Faculty of Science, University of Douala, Cameroon}
\author{Serge Guy Nana Engo}
\affil{Department of Physics, Faculty of Science, University of Yaound\'{e}~I, Cameroon}
\author{Corresponding author: \texttt{steve.teguia@facsciences-uy1.cm}}

\date{\today}

\maketitle

\noindent\textbf{Associated Manuscript:}
This Supporting Information accompanies the manuscript ``Coordinated Vibronic Light-Harvesting and Polaritonic Interface for Symbiotic Quantum Agrivoltaics'' submitted to \textit{Nature Energy}.
Equations, figures, and tables are prefixed with ``S'' (e.g., Eq.~S1, Fig.~S1, Table~S1).
References to the main manuscript are indicated by unadorned equation/figure numbers (e.g., Eq.~1, Fig.~1).

\vspace{1em}

\tableofcontents

\newpage

%==============================================================================
% SECTION S1: 8-SITE FMO HAMILTONIAN AND COUPLING PARAMETERS
%==============================================================================

\section{8-Site FMO Hamiltonian and Full Coupling Matrix}
\label{SI-sec:hamiltonian}

\subsection{Hamiltonian Formulation}

The 8-site Fenna--Matthews--Olson (FMO) Hamiltonian used throughout this work extends the standard 7-site model~\cite{Adolphs2006} by including the eighth bacteriochlorophyll $a$ molecule that bridges the chlorosome baseplate to the reaction center~\cite{Moix2011}.
The full electronic Hamiltonian in the site basis is:

\begin{equation}
\hat{H}_{\mathrm{FMO}} = \sum_{n=1}^{8} \varepsilon_n \dyad{n} + \sum_{n \neq m} J_{nm} \dyad{n}{m} - i\Gamma_{\mathrm{RC}} \sum_{j \in \{3,4\}} \dyad{j}{j},
\label{eq:SI_fmo_hamiltonian}
\end{equation}

where $\varepsilon_n$ are the site energies (referenced to the lowest-energy site), $J_{nm}$ are the inter-site electronic couplings, and $\Gamma_{\mathrm{RC}} = \SI{0.15}{\per\ps}$ is the irreversible reaction center trapping rate applied to Sites~3 and~4.
The anti-Hermitian term $-i\Gamma_{\mathrm{RC}}$ models irreversible charge separation at the reaction center; the Hamiltonian is non-Hermitian only on these diagonal trapping terms.

\subsection{Site Energies}

Table~\ref{tab:SI_site_energies} lists the site energies used in all production simulations.
Energies are given relative to the lowest-energy site (Site~3 at \SI{12210}{\per\cm}) and are expressed in \si{\per\cm}.

\begin{table}[htbp]
\centering
\caption{\textbf{8-site FMO site energies} (in \si{\per\cm}, referenced to Site~3).}
\label{tab:SI_site_energies}
\begin{tabular}{lcccccccc}
\toprule
\textbf{Parameter} & \textbf{Site 1} & \textbf{Site 2} & \textbf{Site 3} & \textbf{Site 4} & \textbf{Site 5} & \textbf{Site 6} & \textbf{Site 7} & \textbf{Site 8} \\
\midrule
Absolute $\varepsilon_n$ & 12410 & 12530 & 12210 & 12320 & 12480 & 12630 & 12440 & 12410 \\
Relative $\varepsilon_n$ & 280   & 420   & 0     & 110   & 270   & 500   & 310   & 200   \\
\bottomrule
\end{tabular}
\end{table}

The absolute site energies are taken from the Adolphs--Renger parameterization~\cite{Adolphs2006} for Sites~1--7.
Site~8, which bridges the baseplate to the reaction center, is assigned an energy of \SI{12410}{\per\cm} based on the eighth-bacteriochlorophyll characterization by Moix~\textit{et al.}~\cite{Moix2011}.

\subsection{Electronic Coupling Matrix}

Table~\ref{tab:SI_coupling_matrix} presents the full $8\times8$ electronic coupling matrix $J_{nm}$ (in \si{\per\cm}).
The matrix is symmetric ($J_{nm} = J_{mn}$), and all 28 unique coupling pairs are listed.
Values for Sites~1--7 follow the Adolphs--Renger parameterization~\cite{Adolphs2006}; couplings involving Site~8 are taken from~Ref.~\cite{Moix2011}.

\begin{table}[htbp]
\centering
\caption{\textbf{Full 8-site electronic coupling matrix $J_{nm}$} (in \si{\per\cm}).}
\label{tab:SI_coupling_matrix}
\begin{tabular}{lcccccccc}
\toprule
 & \textbf{Site 1} & \textbf{Site 2} & \textbf{Site 3} & \textbf{Site 4} & \textbf{Site 5} & \textbf{Site 6} & \textbf{Site 7} & \textbf{Site 8} \\
\midrule
\textbf{Site 1} & 0 & -87.7 & 5.5  & -5.9 & 6.7  & -13.7 & -9.9 & 4.0 \\
\textbf{Site 2} &   & 0     & 30.8 & 8.2  & 0.7  & 11.8  & 4.3  & -2.1 \\
\textbf{Site 3} &   &       & 0    & -53.5 & -2.2 & -9.6  & 6.0  & 12.0 \\
\textbf{Site 4} &   &       &      & 0    & -70.7 & -17.0 & -63.3 & 5.5 \\
\textbf{Site 5} &   &       &      &      & 0    & 81.1  & -1.3 & 3.2 \\
\textbf{Site 6} &   &       &      &      &      & 0     & 39.7 & -8.7 \\
\textbf{Site 7} &   &       &      &      &      &       & 0    & -15.0 \\
\textbf{Site 8} &   &       &      &      &      &       &      & 0 \\
\bottomrule
\end{tabular}
\end{table}

\subsection{Reaction Center Trapping}

The irreversible trapping is implemented through a non-Hermitian term on Sites~3 and~4, which together form the primary funnel to the reaction center in the 8-site FMO model~\cite{Renger2003,Moix2011}.
The trapping rate $\Gamma_{\mathrm{RC}} = \SI{0.15}{\per\ps}$ corresponds to a trapping time of $\tau_{\mathrm{RC}} \approx \SI{6.7}{\ps}$, consistent with ultrafast charge separation rates measured in reaction center complexes~\cite{Renger2003}.
The forward transfer yield is defined as:

\begin{equation}
\Phi_{\mathrm{FT}} = 2\Gamma_{\mathrm{RC}} \int_0^{t_{\mathrm{max}}} \left[ P_3(t) + P_4(t) \right] dt,
\label{eq:SI_phi_ft}
\end{equation}

with $t_{\mathrm{max}} = \SI{1}{\ps}$ and $P_j(t) = \rho_{jj}(t)$.
The factor of 2 accounts for the two trapping sites.

\subsection{Hermiticity Validation}

The non-Hermiticity is strictly confined to the diagonal entries of Sites~3 and~4.
For all $i \neq j$, $H_{ij} = H_{ji}^*$ is satisfied within machine precision (\num{1e-14}).
The Hamiltonian's Hermiticity condition is verified for all 64 matrix elements at every time step of the simulation.

%==============================================================================
% SECTION S2: CONVERGENCE ANALYSIS
%==============================================================================

\section{Convergence Analysis}
\label{SI-sec:convergence}

\subsection{Hierarchy Depth Convergence}

We verify convergence of the PT-HOPS/SBD hierarchy depth $L$ by comparing exciton population dynamics at $L=6, 7, 8$.
Table~\ref{tab:SI_hierarchy_convergence} summarizes the mean absolute error (MAE) between successive $L$ values and the relative deviation from the $L=9$ reference (computed for a subset of 10 disorder realizations).

\begin{table}[htbp]
\centering
\caption{\textbf{Hierarchy depth convergence.} MAE between successive hierarchy depths for exciton populations, averaged over all 8 sites and 100 disorder realizations.}
\label{tab:SI_hierarchy_convergence}
\begin{tabular}{lccc}
\toprule
$L$ & MAE ($L \to L+1$) & Rel.\ dev.\ from $L=9$ & CPU time (min) \\
\midrule
6 & \num{9.8e-10} & \num{1.2e-8} & 2.1 \\
7 & \num{9.6e-10} & \num{1.1e-8} & 4.3 \\
8 & \num{3.1e-11} & \num{4.2e-10} & 7.5 \\
\bottomrule
\end{tabular}
\end{table}

The MAE between $L=7$ and $L=8$ is $\num{3.1e-11}$, which is below double-precision round-off for most observables.
We adopt $L=8$ as the production standard to ensure robust convergence across all disorder realizations.

\subsection{Matsubara Truncation}

We test Matsubara term counts $K = 1, 2, 3$ to assess the impact of high-temperature corrections.
Table~\ref{tab:SI_matsubara_convergence} shows the effect on the coherence lifetime $\tau_c$ and the forward transfer yield $\Phi_{\mathrm{FT}}$.

\begin{table}[htbp]
\centering
\caption{\textbf{Matsubara truncation convergence.} Impact of $K$ on key observables at \SI{295}{\kelvin}.}
\label{tab:SI_matsubara_convergence}
\begin{tabular}{lccc}
\toprule
$K$ & $\tau_c$ (fs) & $\Phi_{\mathrm{FT}}$ & MAE from $K=3$ \\
\midrule
1 & 406(30) & 0.86(3) & \num{3.4e-4} \\
2 & 420(35) & 0.89(3) & \num{3.3e-5} \\
3 & 422(36) & 0.89(3) & --- \\
\bottomrule
\end{tabular}
\end{table}

The MAE between $K=2$ and $K=3$ is $\num{3.3e-5}$, confirming that $K=2$ is sufficient for converged dynamics at room temperature.
This is consistent with the expectation that at \SI{295}{\kelvin}, the dynamics are dominated by the explicit vibronic modes rather than high-order Matsubara corrections to the Drude--Lorentz background.

\subsection{Time Step Convergence}

We verify time step convergence by comparing $\Delta t = \SIlist{0.5;1.0;2.0}{\fs}$.
Table~\ref{tab:SI_timestep_convergence} summarizes the results.

\begin{table}[htbp]
\centering
\caption{\textbf{Time step convergence.} Maximum population deviation relative to the $\Delta t = \SI{0.5}{\fs}$ reference.}
\label{tab:SI_timestep_convergence}
\begin{tabular}{lc}
\toprule
$\Delta t$ (fs) & Max.\ population deviation \\
\midrule
0.5 & --- (reference) \\
1.0 & \num{4.3e-3} \\
2.0 & \num{8.1e-3} \\
\bottomrule
\end{tabular}
\end{table}

We adopt $\Delta t = \SI{0.5}{\fs}$ as the production standard, consistent with Paper~I~\cite{Kouam2026}.

\subsection{Detailed Balance}

The steady-state populations are compared with the analytical Boltzmann distribution to verify detailed balance.
The maximum deviation across all 8 sites is \SI{2.3}{\percent}, confirming that thermal equilibrium is correctly reproduced.

\subsection{Physical Consistency Checks}

All production simulations satisfy the following consistency checks:
\begin{itemize}
\item \textbf{Trace preservation:} $\max_t |\Tr[\rho(t)] - 1| < \num{1e-12}$
\item \textbf{Positivity:} $\min_i \lambda_i(t) > -\num{1e-14}$ for all eigenvalues $\lambda_i$ of $\rho(t)$
\item \textbf{Hermiticity:} $\max_{i,j,t} |\rho_{ij}(t) - \rho_{ji}^*(t)| < \num{1e-14}$
\item \textbf{Population realness:} $\max_{i,t} |\Im[\rho_{ii}(t)]| < \num{1e-14}$
\end{itemize}

%==============================================================================
% SECTION S3: POLARON DRESSING MECHANISM
%==============================================================================

\section{Polaron Dressing Mechanism and Residual Reorganization Energy}
\label{SI-sec:polaron}

\subsection{Canonical (Lang--Firsov) Transformation}

The polaron dressing mechanism that underlies the coherence protection reported in the main text can be rigorously understood through the canonical Lang--Firsov transformation~\cite{Chin2013,Huelga2013}.
The full system-bath Hamiltonian is:

\begin{equation}
\hat{H} = \hat{H}_S + \sum_n \dyad{n} \sum_k g_{n,k}(\hat{b}_{n,k}^\dagger + \hat{b}_{n,k}) + \hat{H}_{\mathrm{bath}},
\label{eq:SI_full_hamiltonian}
\end{equation}

where $g_{n,k}$ are the system-bath coupling constants for mode $k$ at site $n$, and $\hat{b}_{n,k}$ ($\hat{b}_{n,k}^\dagger$) are the bath annihilation (creation) operators.

The unitary transformation $\hat{U} = \exp\!\left[\sum_n \dyad{n} \sum_k \frac{g_{n,k}}{\omega_k}(\hat{b}_{n,k}^\dagger - \hat{b}_{n,k})\right]$ partially diagonalizes the system-bath coupling.
After applying $\hat{U}$, the transformed Hamiltonian becomes:

\begin{equation}
\hat{H}' = \hat{U}^\dagger \hat{H} \hat{U} = \hat{H}_S' + \hat{H}_{\mathrm{bath}} + \hat{H}_{\mathrm{res}},
\label{eq:SI_lf_hamiltonian}
\end{equation}

where $\hat{H}_S'$ contains renormalized excitonic couplings (the polaron-shifted energies) and $\hat{H}_{\mathrm{res}}$ represents the residual system-bath coupling after the transformation.

\subsection{Residual Spectral Density}

The key physical insight is that underdamped vibronic modes satisfying the resonance condition $\Delta E_{nm} \approx \hbar\omega_v \pm J_{nm}$ become partially ``dressed'' by the canonical transformation: they are incorporated into $\hat{H}_S'$ and no longer contribute to dephasing in the transformed frame.
The residual spectral density is:

\begin{equation}
J_{\mathrm{res}}(\omega) = J_{\mathrm{bath}}(\omega) - \sum_{k \in \mathcal{S}} \frac{\lambda_k \omega_k^2 \gamma_k}{(\omega - \omega_k)^2 + \gamma_k^2},
\label{eq:SI_residual_spectral}
\end{equation}

where $\mathcal{S}$ denotes the set of underdamped modes that are vibronically dressed.

\subsection{Quantitative Effect}

For the 12-mode Kleinekath\"{o}fer/Coker parameterization at \SI{295}{\kelvin}, the residual reorganization energy under the two excitation scenarios is:

\begin{align}
\lambda_{\mathrm{res}}^{\mathrm{broad}} &= \lambda_{\mathrm{DL}} + \sum_{k=1}^{12} S_k \omega_k \approx \SI{53}{\per\cm}, \label{eq:SI_lambda_broad} \\
\lambda_{\mathrm{res}}^{\mathrm{filt}} &= \lambda_{\mathrm{DL}} + \sum_{k \notin \mathcal{S}} S_k \omega_k \approx \SI{40}{\per\cm}. \label{eq:SI_lambda_filt}
\end{align}

The \SI{20}{\percent} reduction in $\lambda_{\mathrm{res}}$ translates to a \SI{50}{\percent} coherence lifetime extension (since $\tau_c \propto 1/\lambda_{\mathrm{res}}$ in the high-temperature limit).
The dual-band filter at \SI{750}{\nm} and \SI{820}{\nm} selectively excites those vibronic resonances that are most strongly coupled to the energy-transfer pathway (primarily the \SI{180}{\per\cm} Franck--Condon mode at Sites~3/4 and the \SI{740}{\per\cm} ring deformation mode), maximizing the dressing effect.

\subsection{Connection to PT-HOPS/SBD}

In the PT-HOPS/SBD simulation framework, this protection is captured natively: the bath correlation function is decomposed into Drude--Lorentz and underdamped components, and the hierarchy equations propagate the system-bath entanglement without assuming weak coupling.
The resulting trajectories exhibit the suppressed decoherence predicted by the vibronic-basis analysis when the pulse spectrum matches the underdamped resonances.

%==============================================================================
% SECTION S4: DYNAMIC PHOTO-PROTECTION — FLUENT STARK AND OMIT DETAILS
%==============================================================================

\section{Dynamic Photo-Protection: Floquet Stark and OMIT Behavior}
\label{SI-sec:npq}

\subsection{Floquet Stark Detuning}

Under high solar flux ($> \SI{800}{\W\per\m\squared}$), the OPV dielectric constant changes dynamically, inducing a time-dependent Stark shift.
The driving term in the Hamiltonian is:

\begin{equation}
\hat{H}_{\mathrm{Stark}}(t) = V_0 \cos(\omega_{\mathrm{F}} t) \left( \dyad{1}{1} + \dyad{6}{6} \right),
\label{eq:SI_stark_detuning}
\end{equation}

with $V_0 = \SI{0.05}{\eV}$ and $\omega_{\mathrm{F}} = \SI{1.2}{\THz}$.
At $t = 0$ (zero flux), the cosine term evaluates to $V_0$, ensuring continuity with the undriven Hamiltonian.
When the flux drops below threshold, the amplitude $V_0$ decays to zero with a relaxation time $\tau_{\mathrm{relax}} = \SI{10}{\ps}$, preventing numerical singularities.

\subsection{OMIT Transmission Modulation}

The optomechanically induced transparency (OMIT) mechanism~\cite{Aspelmeyer2014} creates destructive interference in the NPoM cavity that suppresses transmission at the resonant channels.
The transmission coefficient is:

\begin{equation}
T_{\mathrm{OMIT}}(I) = 
\begin{cases}
T_0, & I = 0, \\
T_0 \left(1 + \frac{I}{I_{\mathrm{sat}}}\right)^{-1}, & I > 0,
\end{cases}
\label{eq:SI_omit}
\end{equation}

where $I_{\mathrm{sat}} = \SI{800}{\W\per\m\squared}$.
The piecewise definition at $I=0$ ensures continuity: $T_{\mathrm{OMIT}}(0) = T_0$, corresponding to the undriven cavity transmission.
At high flux (\SI{1000}{\W\per\m\squared}), $T = T_0 \times (1 + 1000/800)^{-1} = 0.44\,T_0$, i.e., transmission drops to \SI{44}{\percent} of baseline.
The time-dependent flux follows $\SI{24}{\hour}$ solar cycles, and the OMIT response is instantaneous on agricultural timescales.

\subsection{Night Mode (Zero Flux)}

At zero flux (night mode), both photo-protection mechanisms return to their baseline states:
\begin{itemize}
\item Floquet Stark: $V_0 \to 0$, restoring the undriven OPV--FMO coupling.
\item OMIT: $T \to T_0$, restoring full transmission.
\end{itemize}
No singularities or discontinuities arise because both models are explicitly defined for $I=0$.

%==============================================================================
% SECTION S5: SOCIOECONOMIC ANALYSIS — COOPERATIVE PAYBACK MODEL
%==============================================================================

\section{Socioeconomic Analysis: Cooperative Payback Model}
\label{SI-sec:socioeconomic}

\subsection{Capital Expenditure and Cost Structure}

The capital expenditure (CAPEX) for the quantum agrivoltaic system (Scenario~A) comprises:
\begin{itemize}
\item \textbf{NPoM-coated graphene barrier:} \SI{120}{\per\m\squared} (including gold nanoparticle synthesis and deposition)
\item \textbf{OPV module:} \SI{80}{\per\m\squared} (blade-coated semi-transparent polymer cells)
\item \textbf{Installation and mounting:} \SI{50}{\per\m\squared}
\item \textbf{Electronics and IoT:} \SI{30}{\per\m\squared} (sensors, microcontroller, QKD module)
\item \textbf{Contingency (\SI{15}{\percent}):} \SI{42}{\per\m\squared}
\end{itemize}
Total CAPEX: \SI{322}{\per\m\squared}.

\subsection{Cooperative Ownership Model}

We model a cooperative of $N$ smallholder farmers who collectively own and operate the quantum agrivoltaic system.
The system size is assumed to scale with cooperative membership: a 5-member cooperative manages a \SI{500}{\m\squared} installation.

The payback period is computed as:

\begin{equation}
P = \frac{C_{\mathrm{total}} (1 - s)}{N \times (R_{\mathrm{elec}} + R_{\mathrm{water}} + R_{\mathrm{crop}})},
\label{eq:SI_payback}
\end{equation}

where $C_{\mathrm{total}} = \SI{322}{\per\m\squared} \times A$ is the total CAPEX for area $A$, $s$ is the subsidy rate (fraction of CAPEX covered by carbon-linked subsidies), $R_{\mathrm{elec}}$ is annual electricity revenue, $R_{\mathrm{water}}$ is water savings value, and $R_{\mathrm{crop}}$ is crop yield improvement revenue.

\subsection{Revenue Streams}

\begin{itemize}
\item \textbf{Electricity:} \SI{180}{\kWh\per\m\squared\per\yr} at \SI{0.12}{\per\kWh} = \SI{21.6}{\per\m\squared\per\yr}
\item \textbf{Water savings:} \SI{1300}{\L\per\m\squared\per\yr} at \SI{0.002}{\per\L} = \SI{2.6}{\per\m\squared\per\yr}
\item \textbf{Crop yield improvement:} \SI{15}{\percent} increase at \SI{1.5}{\per\kg} = \SI{5.0}{\per\m\squared\per\yr}
\item \textbf{Carbon credit:} \SI{12.4}{\kg\ CO2e\per\m\squared\per\yr} at \SI{50}{\per\tonne} = \SI{0.62}{\per\m\squared\per\yr}
\end{itemize}
Total annual revenue: \SI{29.8}{\per\m\squared\per\yr}.

\subsection{Payback Periods}

Table~\ref{tab:SI_payback} summarizes the payback period for various cooperative sizes and subsidy rates.

\begin{table}[htbp]
\centering
\caption{\textbf{Payback period (years) as a function of cooperative size $N$ and subsidy rate $s$.}}
\label{tab:SI_payback}
\begin{tabular}{lcccc}
\toprule
$N$ & $s = \SI{0}{\percent}$ & $s = \SI{20}{\percent}$ & $s = \SI{30}{\percent}$ & $s = \SI{50}{\percent}$ \\
\midrule
1 (individual) & 14.2 & 10.8 & 8.2 & 5.1 \\
3 & 7.8 & 5.9 & 4.5 & 2.8 \\
5 & 5.9 & 4.5 & \textbf{3.5} & 2.1 \\
10 & 4.2 & 3.2 & 2.4 & 1.5 \\
\bottomrule
\end{tabular}
\end{table}

For a 5-member cooperative with a \SI{30}{\percent} carbon-linked subsidy, the payback period of \SI{3.5}{\yr} makes the quantum agrivoltaic system economically viable for smallholder farmers.
This model is conservative: it does not account for potential revenue from the NEB carbon credit market, which could further reduce the payback period to below \SI{3}{\yr}.

\subsection{Tiered Subsidy Framework}

We propose a tiered subsidy framework where the NEB serves as the verifiable metric for subsidy qualification:

\begin{itemize}
\item \textbf{Tier~1 (NEB $<$ \SI{5}{\kg\ CO2e\per\m\squared\per\yr}):} \SI{0}{\percent} subsidy
\item \textbf{Tier~2 (NEB \SIrange{5}{10}{\kg\ CO2e\per\m\squared\per\yr}):} \SI{15}{\percent} subsidy
\item \textbf{Tier~3 (NEB $>$ \SI{10}{\kg\ CO2e\per\m\squared\per\yr}):} \SI{30}{\percent} subsidy
\end{itemize}

This creates a direct incentive loop: higher quantum yield improvements lead to higher NEB, which unlocks greater subsidy access, enabling broader deployment.

%==============================================================================
% SECTION S6: BB84 QKD IMPLEMENTATION DETAILS
%==============================================================================

\section{BB84 Quantum Key Distribution: Implementation and Performance}
\label{SI-sec:qkd}

\subsection{Protocol Details}

The BB84 protocol~\cite{Bennett1984} is implemented as follows:

\begin{enumerate}
\item \textbf{Preparation:} Alice (sensor node) generates $4N = 1024$ random bits (for $N = 256$ final key length) and encodes each bit in one of four polarization states: $|H\rangle$, $|V\rangle$ (rectilinear basis, $+$) or $|D\rangle$, $|A\rangle$ (diagonal basis, $\times$).
\item \textbf{Transmission:} The quantum states are transmitted over a noisy classical channel with a simulated depolarizing noise rate of $p = \SI{5}{\percent}$, modeling the decoherence in field-deployed fiber links~\cite{Zahidy2024}.
\item \textbf{Measurement:} Bob (cloud server) measures each received qubit in a randomly chosen basis ($+$ or $\times$), independent of Alice's preparation basis.
\item \textbf{Sifting:} Alice and Bob publicly announce their basis choices and retain only events where bases match (approximately $2N = 512$ bits remain).
\item \textbf{Parameter estimation:} A random sample (minimum 1 bit, maximum 50 bits) of the sifted key is disclosed to estimate the quantum bit error rate (QBER). The sample fraction follows $f_{\mathrm{sample}} = \min(50/N, 0.2)$.
\item \textbf{Key generation:} If QBER $< \SI{11}{\percent}$ (the security threshold), the remaining sifted bits form the raw key. Error correction and privacy amplification reduce the key to $N = 256$ bits.
\end{enumerate}

\subsection{Error Threshold and Security}

The \SI{11}{\percent} QBER threshold is derived from the intercept-resend attack bound: an eavesdropper (Eve) performing intercept-resend on all qubits introduces a minimum QBER of \SI{25}{\percent}. The \SI{11}{\percent} threshold provides a safety margin for practical implementations~\cite{Bennett1984,Zahidy2024}.

\subsection{Fail-Safe Mechanism}

If the QBER exceeds \SI{11}{\percent} (e.g., due to extreme weather events, equipment malfunction, or attempted eavesdropping), the key generation aborts and a fail-safe local irrigation routine activates:

\begin{itemize}
\item A pre-programmed irrigation schedule (stored in non-volatile memory) ensures uninterrupted water supply to crops.
\item The fail-safe schedule is based on the FAO-56 ET$_c$ estimate for the current season.
\item Telemetry data is cached locally and transmitted when the secure channel is restored.
\end{itemize}

\subsection{Computational Overhead}

Table~\ref{tab:SI_qkd_overhead} shows the runtime overhead of the BB84 simulation for a resource-limited IoT microcontroller (simulated on a single ARM Cortex-M4 core at \SI{120}{\MHz}).

\begin{table}[htbp]
\centering
\caption{\textbf{BB84 QKD simulation runtime overhead.}}
\label{tab:SI_qkd_overhead}
\begin{tabular}{lcc}
\toprule
\textbf{Operation} & \textbf{Time (\si{\milli\second})} & \textbf{Memory (\si{\kilo\byte})} \\
\midrule
State preparation (1024 bits) & 0.8 & 0.5 \\
Noisy transmission & 1.2 & 0.1 \\
Measurement and sifting & 0.6 & 0.3 \\
QBER estimation & 0.1 & 0.1 \\
Key distillation & 2.1 & 1.2 \\
\midrule
Total & 4.8 & 2.2 \\
\bottomrule
\end{tabular}
\end{table}

The total runtime of \SI{4.8}{\ms} and memory footprint of \SI{2.2}{\kB} are negligible for resource-limited IoT nodes, confirming the practical feasibility of BB84-secured agricultural telemetry.

%==============================================================================
% SECTION S7: RESONANCE CONDITION FOR GENERAL VIBRONIC COUPLING
%==============================================================================

\section{Generalized Resonance Condition for Vibronic Coherence Protection}
\label{SI-sec:resonance}

\subsection{Theoretical Derivation}

The resonance condition that enables the coherence protection mechanism is a general property of vibronically coupled chromophore aggregates~\cite{Chin2013,Christensson2012}.
For a dimer of two chromophores with site energies $\varepsilon_1$, $\varepsilon_2$ and electronic coupling $J_{12}$, the excitonic energy gap is:

\begin{equation}
\Delta E_{\mathrm{exc}} = \sqrt{(\varepsilon_1 - \varepsilon_2)^2 + 4J_{12}^2}.
\label{eq:SI_exciton_gap}
\end{equation}

When this gap is resonant with an underdamped vibrational mode of frequency $\omega_v$:

\begin{equation}
\Delta E_{\mathrm{exc}} \approx \hbar\omega_v \pm J_{12},
\label{eq:SI_resonance_condition}
\end{equation}

the excitonic and vibrational states hybridize to form a vibronic (polaron-dressed) state.
The $\pm J_{12}$ term accounts for Davydov splitting in the dimer: the symmetric and antisymmetric combinations of the two site excitations have energies separated by $2J_{12}$, each coupling differently to the vibrational manifold.

In the vibronic basis (after the Lang--Firsov transformation), the effective pure-dephasing rate between excitonic states $\ket{\alpha}$ and $\ket{\beta}$ is:

\begin{equation}
\Gamma_{\alpha\beta}^{\mathrm{pd}} = \frac{1}{\hbar^2} \int_0^\infty d\omega\, \frac{J_{\mathrm{res}}(\omega)}{\omega^2} \left[\coth\!\left(\frac{\beta\hbar\omega}{2}\right)(1 - \cos\omega t) + i\sin\omega t\right]\bigg|_{t\to\infty},
\label{eq:SI_dephasing_rate}
\end{equation}

where $J_{\mathrm{res}}(\omega)$ is the residual spectral density defined in Eq.~\ref{eq:SI_residual_spectral}.
Modes satisfying Eq.~\ref{eq:SI_resonance_condition} are removed from $J_{\mathrm{res}}(\omega)$, reducing $\Gamma_{\alpha\beta}^{\mathrm{pd}}$ proportionally.

\subsection{Generality Across Material Platforms}

The resonance condition is not specific to the FMO complex.
Any molecular aggregate satisfying the condition $\Delta E_{\mathrm{exc}} \approx \hbar\omega_v \pm J_{nm}$ can benefit from spectral coherence protection.
Relevant platforms include:

\begin{itemize}
\item \textbf{Synthetic J-aggregates:} Organic dyes that form excitonically coupled chains with tunable $J_{nm}$~\cite{Jang2008}.
\item \textbf{Plant light-harvesting complexes:} LHCII, CP29, and PSI core complexes that have underdamped vibrational modes in the \SIrange{100}{1600}{\per\cm} range.
\item \textbf{Conjugated polymer domains:} P3HT, PTB7, and other OPV donor polymers where vibronic coupling influences charge separation yields.
\end{itemize}

\subsection{Dual-Band Transmission Design Criterion}

For the OPV spectral filter design, the transmission bands are chosen to satisfy the resonance condition for the dominant energy-transfer pathway in the target photosynthetic complex.
The general design criterion is:

\begin{equation}
\lambda_{\mathrm{filter}} \approx \frac{2\pi c}{\omega_{\mathrm{ex}} \pm \omega_v},
\label{eq:SI_filter_design}
\end{equation}

where $\omega_{\mathrm{ex}}$ is the primary exciton absorption frequency and $\omega_v$ is the underdamped mode frequency.
For FMO at \SI{295}{\kelvin}, $\omega_{\mathrm{ex}} \approx \SI{12500}{\per\cm}$ (\SI{800}{\nm}) and $\omega_v = \SI{180}{\per\cm}$ gives $\lambda_{\mathrm{filter}} \approx \SIlist{790;810}{\nm}$, which defines the dual-band filter passband.

%==============================================================================
% SECTION S8: THREE-SITE EXCITONIC TRIMER VERIFICATION
%==============================================================================

\section{Three-Site Excitonic Trimer Verification}
\label{SI-sec:3site}

\subsection{Motivation}

To confirm that the dual-band filtering effect is not an artifact of the specific FMO complex parameters, we apply the protocol to a minimal three-site excitonic trimer with periodic boundary conditions~\cite{Mohseni2008,Rebentrost2009,Chin2010}.

\subsection{Trimer Parameters}

The trimer Hamiltonian is:

\begin{equation}
\hat{H}_{\mathrm{trimer}} = \sum_{n=1}^{3} \varepsilon_n \dyad{n} + J \sum_{n=1}^{3} \left( \dyad{n}{n+1} + \dyad{n+1}{n} \right),
\label{eq:SI_trimer_hamiltonian}
\end{equation}

with $\varepsilon_n = \varepsilon_0 + \delta_n$ (uniform site energy $\varepsilon_0 = \SI{12410}{\per\cm}$ plus small Gaussian disorder $\delta_n \sim \mathcal{N}(0, \SI{10}{\per\cm})$), and $J = \SI{-50}{\per\cm}$.

The spectral density is a single Drude--Lorentz mode ($\lambda = \SI{20}{\per\cm}$, $\gamma = \SI{30}{\per\cm}$) plus one underdamped mode at $\omega_v = \SI{200}{\per\cm}$ (Huang--Rhys factor $S = 0.08$).

\subsection{Results}

Table~\ref{tab:SI_trimer_results} compares the forward transfer yield enhancement under broadband and filtered excitation.

\begin{table}[htbp]
\centering
\caption{\textbf{Three-site trimer: forward transfer yield under broadband and filtered excitation.}}
\label{tab:SI_trimer_results}
\begin{tabular}{lccc}
\toprule
Excitation & $\Phi_{\mathrm{FT}}$ & $\tau_c$ (fs) & $\eta_{\mathrm{FT}}$ \\
\midrule
Broadband & 0.62(4) & 180(20) & --- \\
Filtered (dual-band) & 0.73(3) & 240(25) & 0.18(3) \\
\bottomrule
\end{tabular}
\end{table}

The coherence enhancement $\eta_{\mathrm{FT}} = (\Phi_{\mathrm{FT}}^{\mathrm{filt}} - \Phi_{\mathrm{FT}}^{\mathrm{broad}}) / \Phi_{\mathrm{FT}}^{\mathrm{broad}} = 0.18(3)$ is consistent with the quantum phase synchronization mechanism identified in photosynthetic antennas~\cite{Zhu2024}.
This demonstrates that the spectral filtering effect is a general vibronic phenomenon, not specific to the FMO complex.

%==============================================================================
% REFERENCES
%==============================================================================

\bibliographystyle{unsrt}
\bibliography{references}

\end{document}
